\section{Contextualização e Problema}
  \label{subsec:contextualizacao-problema}

  A análise de dados é um dos principais fatores para o sucesso de uma empresa, pois é através dela que é possível obter informações que se tornam valiosas para a tomada de decisões. Logo, dado o aumento da quantidade de dados gerados e a necessidade de se obter informações a partir deles, gerou-se o interesse por ferramentas que auxiliem na análise e visualização de dados.

  Um grande conjunto de ferramentas já está disponível online, porém, a maioria delas não é gratuita, performática e nem de código aberto. A ausência de ferramentas deste tipo é um fator notável, visto que nem sempre as empresas possuem recursos ou pessoal para desenvolver ferramentas deste tipo.

  Nesta perspectiva, surge a oportunidade de se desenvolver uma ferramenta que atenda aos critérios citados anteriormente, e que seja facilmente extensível para novas funcionalidades e diferentes tipos e fontes de dados.



